\DocumentMetadata{
  pdfversion=2.0,
  pdfstandard=ua-2,
  lang=en-US,
  tagging=on,
  tagging-setup={math/setup=mathml-SE,tabsorder=structure}
}
\documentclass[11pt,letterpaper]{article}
\usepackage[margin=0.68in,includefoot]{geometry}
\usepackage{newcomputermodern}
\usepackage{amsmath,amscd,mathtools}
\usepackage{tikz,adjustbox}
\providecommand{\square}{\mdlgwhtsquare}
\usepackage[hidelinks]{hyperref}
\hypersetup{
  pdftitle={Numerical Analysis PhD Exam, May 2011},
  pdfauthor={University of Florida Department of Mathematics},
  pdfsubject={Graduate mathematics examination},
  pdfdisplaydoctitle=true,
  bookmarksopen=true
}
\setcounter{secnumdepth}{0}
\setlength{\parindent}{0pt}
\setlength{\parskip}{0.28em}
\setlength{\emergencystretch}{3em}
\setlength{\leftmargini}{2.15em}
\setlength{\leftmarginii}{2.65em}
\setlength{\leftmarginiii}{3em}
\pagestyle{empty}
\raggedbottom
\allowdisplaybreaks
\NewTaggingSocketPlug{block/list/label}{examproblem}{%
  \tagstructbegin{tag=Lbl}\tagstructend%
  \tagstructbegin{tag=\LBody}%
  \tagstructbegin{tag=\ProblemHeadingTag,title-o={\ProblemHeadingText}}%
  \tagmcbegin{tag=\ProblemHeadingTag,actualtext={\ProblemHeadingText}}%
  #2%
  \tagmcend%
  \tagstructend%
}
\newenvironment{examproblems}[2]{%
  \def\ProblemHeadingTag{#2}%
  \AssignTaggingSocketPlug{block/list/label}{examproblem}%
  \begin{enumerate}%
  \setcounter{enumi}{#1}%
  \renewcommand{\labelenumi}{\textbf{\theenumi.}}%
  \setlength{\topsep}{0.35em}%
  \setlength{\partopsep}{0pt}%
  \setlength{\itemsep}{0.48em}%
  \setlength{\parsep}{0.18em}%
}{\end{enumerate}}
\newenvironment{examsubparts}{%
  \AssignTaggingSocketPlug{block/list/label}{default}%
  \begin{enumerate}%
  \setlength{\topsep}{0.18em}%
  \setlength{\partopsep}{0pt}%
  \setlength{\itemsep}{0.16em}%
  \setlength{\parsep}{0.1em}%
}{\end{enumerate}}
\newenvironment{examinstructions}{%
  \par\begingroup\itshape
}{%
  \par\endgroup\vspace{0.18em}
}
\makeatletter
\renewcommand\subsection{\@startsection{subsection}{2}{0pt}%
  {-1.25ex plus -.25ex minus -.1ex}%
  {0.55ex plus .1ex}%
  {\normalfont\large\bfseries}}
\renewcommand{\@maketitle}{%
  \newpage\null\vskip -1em%
  {\footnotesize\bfseries
    \href{https://www.ufl.edu/}{University of Florida}, \href{https://math.ufl.edu/}{Department of Mathematics}\par}%
  \vskip -0.5em%
  \tagstructbegin{tag=H1,title={Numerical Analysis PhD Exam, May 2011}}%
  \tagpdfparaOff%
  \tagmcbegin{tag=H1}%
    {\LARGE\bfseries\href{https://gma.math.ufl.edu/exams/numerical-analysis-phd/}{Numerical Analysis PhD Exam}, May 2011\par}%
  \tagmcend%
  \tagpdfparaOn%
  \tagstructend%
  \vskip 0.25em%
  \tagmcbegin{artifact}%
    \hrule height 1pt%
  \tagmcend%
  \vskip 1em%
}
\makeatother
\title{Numerical Analysis PhD Exam, May 2011}
\author{}
\date{}
\begin{document}
\maketitle
\thispagestyle{empty}
\begin{examinstructions}
Answer any 8 questions.
\end{examinstructions}
\begin{examproblems}{0}{H2}
\def\ProblemHeadingText{Problem 1.}
\item
Suppose~\(A\) and~\(B\) are square matrices such that \(AB\) is normal. Prove that \(\|AB\|_2 \leq \|BA\|\). (We use \(\|\cdot\|_2\) to denote the spectral norm and \(\|\cdot\|\) to denote any induced matrix norm.)
\def\ProblemHeadingText{Problem 2.}
\item
Let~\(P\) and~\(Q\) be Hermitian positive definite matrices. Prove that
\[
x^*Px \leq x^*Qx \quad \text{for all } x \in \mathbb{C}^n
\]
 if and only if
\[
x^*Q^{-1}x \leq x^*P^{-1}x \quad \text{for all } x \in \mathbb{C}^n.
\]
\def\ProblemHeadingText{Problem 3.}
\item
Suppose~\(A\) is a Hermitian positive definite matrix split into \(A=C+C^*+D\) where~\(D\) is also Hermitian positive definite. Prove that \(B=C+\omega^{-1}D\) is invertible whenever \(0<\omega<2\). Consider the iteration \(x_{n+1}=x_n+B^{-1}(b-Ax_n)\), with any initial iterate \(x_0\). Prove that \(x_n\) converges to \(x=A^{-1}b\) whenever \(0<\omega<2\).
\def\ProblemHeadingText{Problem 4.}
\item
Let the singular values of any matrix \(M\in\mathbb{C}^{m\times n}\) be denoted by \(\sigma_1(M)\geq\sigma_2(M)\geq\cdots\geq\sigma_q(M)\) with \(q=\min(m,n)\). Prove that if~\(A\) and~\(B\) are two matrices in \(\mathbb{C}^{m\times n}\), then
\[
\sigma_{i+j-1}(A+B)\leq\sigma_i(A)+\sigma_j(B)
\]
 for all \(i,j=1,2,\ldots,q\) and \(i+j\leq q\).
\def\ProblemHeadingText{Problem 5.}
\item
Let \(A\in\mathbb{R}^{N\times N}\) and \(b\in\mathbb{R}^N\). Consider the following iteration for solving \(Ax=b\), that computes \(x_{n+1}\), given \(x_n\in\mathbb{R}^N\), as follows (in~\(m\) intermediate steps): Setting \(x_{n+1}^{(0)}=x_n\), for \(\ell=1,2,\ldots,m\), compute \(x_{n+1}^{(\ell)}=x_{n+1}^{(\ell-1)}+\tau_\ell(b-Ax_{n+1}^{(\ell-1)})\). Then, define \(x_{n+1}=x_{n+1}^{(m)}\). There is a linear operator~\(E\) such that \(x_{n+1}-x=E(x_n-x)\). Give a formula for~\(E\). Suppose~\(A\) is Hermitian and positive definite with spectral condition number~\(\kappa\). Prove that there are real values of the~\(m\) parameters \(\tau_\ell\) such that
\[
\rho(E)\leq 2\left(\frac{\sqrt{\kappa}-1}{\sqrt{\kappa}+1}\right)^m.
\]
\def\ProblemHeadingText{Problem 6.}
\item
Let \(I_n(f)=\sum_{j=0}^n w_{j,n}f(x_{j,n})\) be a sequence of quadratures on \([a,b]\) such that (i) \(I_n(p)\to I(p)\) as \(n\to+\infty\) for any polynomial \(p(x)\), and (ii) \(B=\sup_n\sum_{j=0}^n|w_{j,n}|<+\infty\). Prove that \(I_n(f)\to I(f)\) for all \(f\in C[a,b]\). Here, the notation \(I(f)=\int_a^b f(x)\,dx\) is used.
\def\ProblemHeadingText{Problem 7.}
\item
Design a stable quadratic algorithm for computing the positive root~\(\alpha\) of the equation \(x^2+2px-q=0\) where~\(p\) and~\(q\) are arbitrary positive numbers. Prove that your algorithm is indeed stable and quadratic. Find an interval \([a,b]\) such that any iterative sequence starting in \([a,b]\) will converge to~\(\alpha\).
\def\ProblemHeadingText{Problem 8.}
\item
State the Lagrange interpolation problem. Prove that the problem is well-posed, i.e. prove the existence and uniqueness of solutions. Derive the error formula for the interpolating polynomial.
\def\ProblemHeadingText{Problem 9.}
\item
Describe the Trapezoidal Rule for numerical integration. State and prove the error formula of the Trapezoidal Rule. Show that the error bound cannot be improved.
\def\ProblemHeadingText{Problem 10.}
\item
Let \(w(x)>0\) be integrable on \([a,b]\). Define an orthogonal polynomial family \(\{\phi_n\}\) on \([a,b]\) with weight \(w(x)\). Explain how \(\{\phi_n\}\) can be constructed from \(\{x^n\}\). Prove that \(\phi_n(x)\) has~\(n\) distinct roots in the interval \((a,b)\).
\end{examproblems}
\end{document}
